%============================================================
% difficulty.tex — coloured difficulty badges and problem links
% used in the pattern EXERCISE lists (E1.x, E2.x, ...).
%
%   \exeasy  \exmedium  \exhard      coloured difficulty badges
%   \exprob{title}{source}{url}      clickable title + source + Solve button
%   \exown{title}                    exercise without an online judge
%============================================================
\definecolor{diffeasy}{RGB}{46, 160, 67}
\definecolor{diffmedium}{RGB}{232, 140, 0}
\definecolor{diffhard}{RGB}{211, 47, 47}

% small coloured pill
\newcommand{\exbadge}[2]{\tikz[baseline=(B.base)]\node[fill=#1, text=white, rounded corners=2pt,
    inner xsep=3pt, inner ysep=1.5pt, font=\scriptsize\bfseries] (B) {\strut #2};}

\definecolor{diffexpert}{RGB}{123, 31, 162}

% The ONE difficulty scale of the book (replaces the old star ratings):
%   \exeasy    green   warm-up, direct use of the pattern
%   \exmedium  orange  needs one idea beyond the template
%   \exhard    red     combines patterns / non-obvious state
%   \exexpert  purple  contest-final / 2400+ level
\newcommand{\exeasy}{\exbadge{diffeasy}{Easy}}
\newcommand{\exmedium}{\exbadge{diffmedium}{Medium}}
\newcommand{\exhard}{\exbadge{diffhard}{Hard}}
\newcommand{\exexpert}{\exbadge{diffexpert}{Expert}}

% Codeforces rating badge in the official Codeforces rank colours:
%   <1200 gray | <1400 green | <1600 cyan | <1900 blue | <2100 violet | <2400 orange | >=2400 red
\definecolor{cfgray}{RGB}{128, 128, 128}
\definecolor{cfgreen}{RGB}{0, 128, 0}
\definecolor{cfcyan}{RGB}{3, 168, 158}
\definecolor{cfblue}{RGB}{0, 0, 255}
\definecolor{cfviolet}{RGB}{170, 0, 170}
\definecolor{cforange}{RGB}{255, 140, 0}
\definecolor{cfred}{RGB}{255, 0, 0}
\newcommand{\cf}[1]{%
    \ifnum#1<1200 \exbadge{cfgray}{CF #1}%
    \else\ifnum#1<1400 \exbadge{cfgreen}{CF #1}%
    \else\ifnum#1<1600 \exbadge{cfcyan}{CF #1}%
    \else\ifnum#1<1900 \exbadge{cfblue}{CF #1}%
    \else\ifnum#1<2100 \exbadge{cfviolet}{CF #1}%
    \else\ifnum#1<2400 \exbadge{cforange}{CF #1}%
    \else \exbadge{cfred}{CF #1}%
    \fi\fi\fi\fi\fi\fi}

% LeetCode's OFFICIAL difficulty, in LeetCode's colours, shown next to the book badge
% (the book badge says where a problem sits in THIS chapter; the LC badge says what LeetCode calls it)
\definecolor{lceasy}{RGB}{0, 184, 163}
\definecolor{lcmedium}{RGB}{255, 192, 30}
\definecolor{lchard}{RGB}{255, 55, 95}
\newcommand{\lc}[1]{%
    \ifstrequal{#1}{Easy}{\exbadge{lceasy}{LC Easy}}{%
    \ifstrequal{#1}{Medium}{\tikz[baseline=(B.base)]\node[fill=lcmedium, text=black, rounded corners=2pt,
        inner xsep=3pt, inner ysep=1.5pt, font=\scriptsize\bfseries] (B) {\strut LC Medium};}{%
    \exbadge{lchard}{LC Hard}}}}

\newcommand{\exprob}[3]{\phantomsection\label{pb:#3}\href{#3}{\textbf{\textcolor{deepblue}{#1}}}
    {\small\textcolor{darkgray}{(#2)}}~\href{#3}{\exbadge{primaryblue}{Solve $\nearrow$}}}
\newcommand{\exown}[1]{\textbf{#1} {\small\textcolor{darkgray}{(own exercise)}}}

% Stretch: a practice problem ABOVE the chapter's level (kept for strong readers; never needed to continue)
\definecolor{diffstretch}{RGB}{0, 105, 120}
\newcommand{\exstretch}{\exbadge{diffstretch}{Stretch}}

% Problem KIND markers. They print NOTHING: they only annotate the source so that tools/check_problem_mix.py
% can check the mix. Every practice set mixes all four kinds:
%   \pclassic  a well-known problem every array programmer should have solved
%   \pconcept  isolates one idea of the chapter in its purest form
%   \precent   from a recent contest (last ~3 years): the current style of problems
%   \ptricky   a trap: an edge case, a wrong-but-tempting idea, or a twist in the model
\newcommand{\pclassic}{\ignorespaces}
\newcommand{\pconcept}{\ignorespaces}
\newcommand{\precent}{\ignorespaces}
\newcommand{\ptricky}{\ignorespaces}

% one-line legend placed at the top of every exercise list
\newcommand{\exlegend}{%
    {\small Difficulty:\; \exeasy~warm-up \quad \exmedium~standard \quad \exhard~challenge \quad \exexpert~expert;
    Codeforces problems also show their rating, e.g.\ \cf{1500}; LeetCode problems show LeetCode's own label, e.g.\ \lc{Medium}.
    Click a problem name or \exbadge{primaryblue}{Solve $\nearrow$} to open it on the judge.}\par\medskip}
